% !TEX root = scientific-computing.tex
%
\hypertarget{ch:packages}{%
\chapter{Loading Packages}\label{ch:packages}}


As with many software packages, additional functionality is available in
other libraries or packages. If you take a look at
\href{Julia's\%20Package\%20Listing}{https://juliaobserver.com} there
are about 2000 packages (as of Fall 2018). Clicking on any of the
package names will give additional information about the package
including documentation (hopefully).

If you don't have a package installed and try to use it, like:

\begin{jlverbatim}
using Primes
\end{jlverbatim}
%
you will get an error that it isn't installed--unless you actually installed it previously. You can add the package by

\begin{jlverbatim}
using Pkg
Pkg.add("Primes")
\end{jlverbatim}

and after fetching the package, rerunning
%
\begin{jlblock}
using Primes
\end{jlblock}
%
should no longer give an error. Then, for example, try
%
\begin{jlblock}
isprime(17)
\end{jlblock}
%
which returns \verb!true! because 17 is a prime number. This is a
nondeterministic function that determines if a number is true. The
documentation for the Primes package is at
\url{https://github.com/JuliaMath/Primes.jl} and demonstrates 4 commands
including \texttt{isprime} and \texttt{factor} which returns all prime
factors (and the order).

A few other packages that we have seen or will see:

\begin{itemize}

\item
  IJulia -- provides the browser interface to Julia. You are probably
  already using this.
\item
  ForwardDiff -- does Automatic Differentiation (see Chapter 8)
\item
  Plots -- A plotting packages (See Chapter 11)
\item
  GR -- A backend to the Plots package.
\item
  PlotlyJS -- Another nice backend to the Plots package.
\item
  DataFrames -- a way to nicely handle datasets generally loaded from
  external files.
\end{itemize}

\href{ch14.html}{Chapter 14} will show how to create our own packages
also called a module.

\hypertarget{managing-packages-in-julia}{%
\section{Managing Packages in
Julia}\label{managing-packages-in-julia}}

Although we can use the \texttt{Pkg} package to handle packages, this
section will use the terminal REPL to handle any non-standard packages.
Additional documentation on this is given in
\href{Julia's\%20package\%20manager\%20help\%20pages}{https://docs.julialang.org/en/latest/stdlib/Pkg/}
First, open up a terminal version of julia (generally by opening the
application that you downloaded). You will get:

\begin{verbatim}
              _
   _      _ _(_)_      |  Documentation: https://docs.julialang.org
  (_)    | (_) (_)     |
  _ _   _| |_  __ _    |  Type "?" for help, "]?" for Pkg help.
  | | | | | | |/ _` |  |
  | | |_| | | | (_| |  |  Version 1.0.0 (2018-08-08)
 _/ |\__'_|_|_|\__'_|  |  Official https://julialang.org/ release
|__/                   |
\end{verbatim}

or similar and then

\begin{verbatim}
julia>
\end{verbatim}

which means we're ready to handle julia commands. If we type
\texttt{{]}}, then the prompt turns into:

\begin{verbatim}
(v1.4) pkg>
\end{verbatim}
% 
where the 1.4 will be the version of julia that you are running. There are a number of commands that we will cover here:

\begin{itemize}
\item  add
\item  remove (rm)
\item  develop (dev)
\item  status
\item  update (up)
\item  test
\item  build
\item  precompile
\end{itemize}
and the commands in parentheses are the shortcut. 

\hypertarget{adding-a-package}{%
\subsection{Adding a package}\label{adding-a-package}}

In the package command line, type \texttt{add} \emph{package\_name} to
add the package. For example, to add the \texttt{ForwardDff} package:

\begin{verbatim}
add ForwardDiff
\end{verbatim}

and it is not installed, you will get something like:

\begin{verbatim}
Updating registry at `~/.julia/registries/General`
  Updating git-repo `https://github.com/JuliaRegistries/General.git`
  Updating `~/.julia/environments/v1.0/Project.toml`
  [f6369f11] + ForwardDiff v0.9.0
  Updating `~/.julia/environments/v1.0/Manifest.toml`
  [9e28174c] + BinDeps v0.8.10
  [bbf7d656] + CommonSubexpressions v0.2.0
  [163ba53b] + DiffResults v0.0.3
  [b552c78f] + DiffRules v0.0.7
  [f6369f11] + ForwardDiff v0.9.0
  [77ba4419] + NaNMath v0.3.2
  [276daf66] + SpecialFunctions v0.7.0
  [90137ffa] + StaticArrays v0.8.3
\end{verbatim}

A few things to note:

\begin{itemize}
\item You results will vary depending on version numbers avaiable and what subpackages (like \verb!DiffRules! or \verb!StaticArrays! are needed for the current version of the package you want to load. 
\item  The line after the 2nd \texttt{Updating} line is the package (and version) that you are installing.
\item  All of the lines after the 3rd \texttt{Updating} line is all of the  packages that this depends on.
\item  The + sign means that the package is being added.
\end{itemize}

If you want to add multiple packages at the same time, say packages A, B
and C, type \texttt{add\ A\ B\ C}. You can also add particular versions
of a package (often for testing or to avoid a bug). For example, if you
want version 0.3.0 of \texttt{ForwardDiff} type:

\begin{verbatim}
add ForwardDiff@0.3.0
\end{verbatim}

You will then get info on the dependencies on that version.

\hypertarget{package-status}{%
\subsection{Package Status}\label{package-status}}

The \texttt{status} command (or \texttt{st}) will just list all of the
main packages installed. For example,

\begin{verbatim}
Status `~/.julia/environments/v1.0/Project.toml`
[c52e3926] Atom v0.7.6
[336ed68f] CSV v0.3.1
[a93c6f00] DataFrames v0.13.1
[f6369f11] ForwardDiff v0.9.0
[7073ff75] IJulia v1.12.0
[e5e0dc1b] Juno v0.5.3
[1a8c2f83] Query v0.10.0
\end{verbatim}

and note that these are just the packages added by the \texttt{add}
command, not all of the dependencies. If you want all of the
dependencies as well, type \texttt{st\ -\/-manifest} and I get a huge
list of packages.

\hypertarget{removing-a-package}{%
\subsection{Removing a Package}\label{removing-a-package}}

You can remove a package by typing \texttt{remove} or \texttt{rm} then
the package name. If I want to remove the \texttt{ForwardDiff} package,
then

\begin{verbatim}
remove ForwardDiff
\end{verbatim}

we get the following:

\begin{verbatim}
Updating `~/.julia/environments/v1.0/Project.toml`
 [f6369f11] - ForwardDiff v0.3.0
 Updating `~/.julia/environments/v1.0/Manifest.toml`
 [49dc2e85] - Calculus v0.4.1
 [c5cfe0b6] - DiffBase v0.2.0
 [f6369f11] - ForwardDiff v0.3.0
 [77ba4419] - NaNMath v0.3.2
\end{verbatim}

Note:

\begin{itemize}

\item
  The \texttt{rm} command removes the package from the list of available
  packages, but doesn't remove them from your harddrive.
\item
  If you want to see everything installed, navigate to the
  \texttt{\textasciitilde{}/.julia/packages} directory, which is where
  they are stored.
\end{itemize}

\hypertarget{updating-packages}{%
\subsection{Updating packages}\label{updating-packages}}

If you type \texttt{update} or \texttt{up} you will update all of the
installed packages (and dependencies). For example:

\begin{verbatim}
Updating `~/.julia/environments/v1.0/Project.toml`
 [7073ff75] ↑ IJulia v1.11.1 ⇒ v1.12.0
 Updating `~/.julia/environments/v1.0/Manifest.toml`
 [7073ff75] ↑ IJulia v1.11.1 ⇒ v1.12.0
 [b85f4697] ↑ SoftGlobalScope v1.0.5 ⇒ v1.0.7
 [5e66a065] ↑ TableShowUtils v0.1.1 ⇒ v0.2.0
\end{verbatim}

and all updates will be with an ↑. If you only want to update a single
package, type the name after \texttt{update}.

\hypertarget{building-packages}{%
\subsection{Building Packages}\label{building-packages}}

Generally a package is built after it is installed. Building a package
might include running code (or unpacking files) after it is installed.
Sometimes if things get wonky, rebuilding is a good thing to do.

\begin{verbatim}
build
\end{verbatim}

or if you only want say \texttt{IJulia} built,

\begin{verbatim}
build IJulia
\end{verbatim}

\subsection{Precompiling packages}

When a package is used, often it requests to be compiled. For example,
when \texttt{using\ Primes}, then following is shown:

\begin{verbatim}
[ Info: Precompiling Primes [27ebfcd6-29c5-5fa9-bf4b-fb8fc14df3ae]
\end{verbatim}

and basically some code is compiled beforehand, generally to speed up
code. You can precompile all code with

\begin{verbatim}
precompile
\end{verbatim}

and it may take a while, but you won't have to wait, when you load the
package with the \texttt{using} command.

\hypertarget{testing-a-packages}{%
\subsection{Testing a Packages}\label{testing-a-packages}}

To test a packge, say the \texttt{ForwardDiff} package, then

\begin{verbatim}
test ForwardDiff
\end{verbatim}

It list all of the dependencies first, and then runs a number of tests
(and we will show how to write tests soon) and timing information. After
a while, it finishes sucessfully.
